\SourceLeaf{243}%
flexible, inextensible solid with one dimension taken to vanish, the properties of the surface depend in part on the shape we suppose it has been brought into. In part they are absolute and remain unchanged whatever shape the surface is bent into. The study of these latter properties opens up a new and fertile field for geometry. They include the measure of curvature and integral curvature in the sense in which we use these expressions. The theory of shortest lines also belongs here, along with several other topics that we reserve for later treatment. From this point of view a plane surface and a surface that can be developed onto a plane, for example a cylindrical surface, a conical surface or the like, are regarded as essentially identical. The proper way of expressing in general the nature of a surface considered in this way always rests on the formula \(\BareRoot{(E\mathrm{d}p^2+2F\mathrm{d}p\mathord{.}\mathrm{d}q+G\mathrm{d}q^2)}\), which relates the element to the two indeterminates \(p,q\). Before pursuing this subject further, however, we must first set out the principles of the theory of shortest lines on a given curved surface.
\Article{14}
A curved line in space is generally specified by expressing the coordinates \(x,y,z\) of each of its points as functions of a single variable, which we shall denote by \(w\). The length of such a line from an arbitrary starting point to the point with coordinates \(x,y,z\) is expressed by the integral
\[
\int\mathrm{d}w\mathord{.}\BareRoot{((\tfrac{\mathrm{d}x}{\mathrm{d}w})^2+(\tfrac{\mathrm{d}y}{\mathrm{d}w})^2+(\tfrac{\mathrm{d}z}{\mathrm{d}w})^2)}
\]
If we suppose that the position of the curved line undergoes an infinitesimal variation, so that the coordinates of each point acquire the variations \(\delta x,\delta y,\delta z\), we find the variation of the total length
\[
=\int\frac{\mathrm{d}x\mathord{.}\mathrm{d}\delta x+\mathrm{d}y\mathord{.}\mathrm{d}\delta y+\mathrm{d}z\mathord{.}\mathrm{d}\delta z}{\BareRoot{(\mathrm{d}x^2+\mathrm{d}y^2+\mathrm{d}z^2)}}
\]
We put this expression in the following form:
\[
\begin{aligned}
&\frac{\mathrm{d}x\mathord{.}\delta x+\mathrm{d}y\mathord{.}\delta y+\mathrm{d}z\mathord{.}\delta z}{\BareRoot{(\mathrm{d}x^2+\mathrm{d}y^2+\mathrm{d}z^2)}}\\
&-\int(\delta x\mathord{.}\mathrm{d}\frac{\mathrm{d}x}{\BareRoot{(\mathrm{d}x^2+\mathrm{d}y^2+\mathrm{d}z^2)}}+\delta y\mathord{.}\mathrm{d}\frac{\mathrm{d}y}{\BareRoot{(\mathrm{d}x^2+\mathrm{d}y^2+\mathrm{d}z^2)}}\\
&\quad+\delta z\mathord{.}\mathrm{d}\frac{\mathrm{d}z}{\BareRoot{(\mathrm{d}x^2+\mathrm{d}y^2+\mathrm{d}z^2)}})
\end{aligned}
\]
When the line is shortest between its endpoints, it is clear that what stands here under the integral sign must vanish. Insofar as the line must lie on a
